% RACER-GT methodology manuscript (English).
% Build:  python scripts/build_docs.py --only methodology --lang en
\documentclass[12pt,a4paper,oneside]{article}
% English setup. Compiled with XeLaTeX so that both languages share one engine
% and one font family, which keeps the two PDFs typographically comparable.
\usepackage{fontspec}
\setmainfont{Noto Serif}
\setsansfont{Noto Sans}
\setmonofont{Noto Sans Mono}

% Author-year citations. Loaded here rather than in _common because it must
% precede hyperref, and because the Chinese manuscript uses numeric citations
% that natbib's author-year mode would reject. Bibliographies are hand-written
% thebibliography blocks: biblatex/biber is absent from a BasicTeX install, and
% a document set that only builds on a full TeX Live is not reproducible enough
% for this project.
\usepackage[authoryear,round]{natbib}

% Single source of truth for version and date across every RACER-GT document.
% Regenerate nothing by hand: bump here, rebuild, and every PDF agrees.
\newcommand{\racerversion}{1.1.0}
\newcommand{\racerdate}{2026-07-27}
\newcommand{\racerauthor}{Yi-Hao Lai}
\newcommand{\raceraffil}{Department of Finance, Da-Yeh University}
\newcommand{\raceremail}{yhlai@mail.dyu.edu.tw}
\newcommand{\racerrepo}{https://github.com/yhlai0911/RACER-GT}

% Shared package set and mathematical macros for every RACER-GT document,
% in both languages. Language-specific font and theorem-name setup lives in
% _preamble-zh.tex and _preamble-en.tex.
%
% Font policy: the build depends only on the five Noto families below. They are
% installable on Linux CI (fonts-noto-core, fonts-noto-cjk) and on macOS via
% Homebrew casks. Do not introduce a font outside this set without adding it to
% .github/workflows/docs.yml as well, or the PDF build stops being reproducible.

\usepackage[margin=2.35cm,headheight=15pt]{geometry}
\usepackage{amsmath,amssymb,amsthm,mathtools,bm}
\usepackage{booktabs,longtable,tabularx,array,multirow}
\usepackage{graphicx,float,caption,subcaption}
\usepackage{enumitem}
\usepackage{xcolor}
\usepackage{microtype}
\usepackage{setspace}
\usepackage{fancyhdr}
\usepackage{hyperref}
\usepackage[nameinlink,noabbrev]{cleveref}
\usepackage{algorithm}
\usepackage{algpseudocode}
\usepackage{listings}
\usepackage{tikz}
\usetikzlibrary{arrows.meta,positioning,shapes.geometric,fit,calc}

\hypersetup{
  colorlinks=true,
  linkcolor=black,
  citecolor=black,
  urlcolor=blue,
  pdfauthor={\racerauthor}
}

\pagestyle{fancy}
\fancyhf{}
\rhead{v\racerversion}
\cfoot{\thepage}
\setstretch{1.28}
\setlength{\parskip}{0.35em}
\setlist[itemize]{leftmargin=2em,itemsep=0.25em,topsep=0.3em}
\setlist[enumerate]{leftmargin=2.2em,itemsep=0.25em,topsep=0.3em}

% ---------------------------------------------------------------- mathematics
\newcommand{\E}{\mathbb{E}}
\newcommand{\Var}{\operatorname{Var}}
\newcommand{\Cov}{\operatorname{Cov}}
\newcommand{\tr}{\operatorname{tr}}
\newcommand{\diag}{\operatorname{diag}}
\newcommand{\argmin}{\operatorname*{arg\,min}}
\newcommand{\argmax}{\operatorname*{arg\,max}}
\newcommand{\one}{\bm{1}}
\newcommand{\Rset}{\mathbb{R}}
\newcommand{\R}{\mathbb{R}}
\newcommand{\plim}{\operatorname*{plim}}
\newcommand{\RACER}{\textsc{Racer-GT}}
\newcommand{\indic}[1]{\mathbb{I}\!\left\{#1\right\}}
\newcommand{\Sig}{\bm{\Sigma}}
\newcommand{\wvec}{\bm{w}}
\newcommand{\evec}{\bm{e}}
\newcommand{\Zvec}{\bm{Z}}

% Design facets, used identically in both languages so that a reader can move
% between the Chinese and English documents without re-learning notation.
\newcommand{\facetT}{\mathcal{T}}   % historical date
\newcommand{\facetD}{\mathcal{D}}   % collection day
\newcommand{\facetS}{\mathcal{S}}   % collection stream
\newcommand{\facetR}{\mathcal{R}}   % technical replicate

\lstset{
  basicstyle=\ttfamily\footnotesize,
  breaklines=true,
  frame=single,
  rulecolor=\color{gray!40},
  backgroundcolor=\color{gray!5},
  columns=fullflexible,
  keepspaces=true,
  showstringspaces=false
}


\setlength{\parindent}{1.5em}

\newtheorem{assumption}{Assumption}[section]
\newtheorem{proposition}{Proposition}[section]
\newtheorem{corollary}{Corollary}[section]
\newtheorem{remark}{Remark}[section]
\newtheorem{definition}{Definition}[section]
\newtheorem{lemma}{Lemma}[section]

\lhead{RACER-GT: Methodology}

\title{\bfseries
RACER-GT: A Design-Based Statistical Framework for the Reliable\\
Construction of Long-Horizon Google Trends Time Series\\[0.6em]
\large Randomized Acquisition, Global Overlap Calibration, Generalizability\\
Decomposition, and Covariance-Adjusted Consensus
}
\author{\racerauthor\\
\small \raceraffil\\
\small \texttt{\raceremail}}
\date{\racerdate\\[0.4em]\small Methodology and software report, version \racerversion}

\begin{document}
\maketitle

\begin{abstract}
\noindent Google Trends (GT) is widely used in finance and economics as a
high-frequency proxy for investor attention, information demand, and sentiment.
Yet the web interface rescales every request to a $0$--$100$ index, a long daily
history must be assembled from several short query windows, and repeating an
identical query can return different numbers. A researcher who downloads a
series once, stitches the windows sequentially, and feeds the result into a
financial model therefore inherits four distinct problems at once: segment-scale
error, dependence across repeated retrievals, zero inflation and rounding at low
search volume, and attenuation bias in any downstream regression.

This paper develops \RACER{} (Randomized Acquisition, Calibration, Error
Decomposition, and Reliability for Google Trends). The framework does not claim
to recover Google's absolute search counts. It instead defines two estimands
that \emph{are} identified---the expected GT return under a fixed query and a
fixed retrieval design, and the latent relative search signal on a common
scale---and then estimates them through eight preregistered stages: (i) a
balanced, protocol-locked repeated-retrieval experiment; (ii) global weighted
least squares on the overlap graph, which uses every usable overlap
simultaneously instead of accumulating error through sequential stitching;
(iii) exhaustive retention of exact duplicates together with a residual-based
near-duplicate dependence graph; (iv) a three-facet Generalizability Theory
decomposition over historical date, collection day, and collection stream;
(v) a design-based stratified reference estimator and a covariance-adjusted
minimum-variance consensus; (vi) quadratic-programming temporal benchmarking
against weekly or monthly aggregates; (vii) separate detection, level, and
daily-innovation reliability with a locked acceptance decision tree; and
(viii) propagation of the resulting measurement uncertainty into
reliability-corrected OLS and SIMEX.

We state the identification conditions explicitly and prove conditional
unbiasedness and minimum-variance results under them. Twenty controlled Monte
Carlo replications, evaluated at matched information sets, give a mixed result.
Cross-pull aggregation cuts mean RMSE against a known latent truth from $7.3519$
for a single pull to $3.7421$, a $49.1\%$ reduction that improves in all twenty
replications, and a further $10.3\%$ below the cross-pull median's $4.1710$.
Temporal benchmarking contributes about $0.22$ more, holding the estimator fixed
($p\approx3\times10^{-12}$). Covariance-adjusted weighting, however, shows no
detectable advantage over the simple cross-pull mean of $3.7253$ ($-0.0167$,
$p=0.152$). The natural explanation is that dependence is too homogeneous for
GLS to exploit, and we test it directly in a four-scenario experiment that puts
a subset of pulls behind a shared cache disturbance cutting across the design
facets. The explanation \textbf{fails}: GLS wins no heterogeneous scenario
significantly and is significantly worse in two. Diagnostics show the
constraints are not responsible (the unrestricted solution is identical to the
digit) and that the mechanism works as intended (weights correlate $-0.562$ with
shared residual dependence), leaving the error in estimating $\Sig$ from nine
pulls as what exceeds the efficiency it buys. The theory is untouched --- GLS is
still BLUE with $\Sig$ known --- but the practical recommendation is not. The
calibration stage yields the same shape of result: global overlap-graph
calibration does stop error accumulating along the chain of joins relative to
sequential stitching (growth ratio $1.12$ against $4.75$), yet the difference in
reconstruction error reaches only $p=0.098$ under the setting most favourable to
it. Taken together these give one cross-cutting conclusion: each mechanism works
as designed and can be shown to work, while its effect on point accuracy is
second order. The value of this framework is that it produces a series whose
quality can be \textbf{judged}, not a more accurate one. None of these figures is
a universal optimality guarantee across keywords, regions, or GT versions.

\medskip
\noindent\textbf{Keywords:} Google Trends; investor attention; measurement
error; experimental design; generalizability theory; graph calibration;
generalized least squares; temporal benchmarking; SIMEX.

\medskip
\noindent\textbf{JEL classification:} C18, C43, C81, G14.
\end{abstract}

\tableofcontents
\newpage

% =============================================================================
\section{Introduction}

Search data have become a standard empirical input in finance. \citet{DaEG2011,DaEG2015}
showed that Google search frequency for ticker symbols captures retail investor
attention in a way that pre-existing proxies do not; \citet{PreisMS2013}
documented associations between search volume for finance-related terms and
subsequent market movements; \citet{ChoiVarian2012} established the broader
``predicting the present'' agenda. A large applied literature now treats a
downloaded GT series as if it were an observed variable measured without error.

That treatment is difficult to defend once the retrieval mechanism is examined.
Three features of the GT interface interact badly with long-horizon daily
research designs.

\paragraph{Per-request renormalization.} Every GT response is rescaled so that
the maximum value within the requested window equals $100$. The returned numbers
are therefore not comparable across requests with different windows. Two
requests that differ only in their end date produce two different scales for the
same calendar day.

\paragraph{Window-length quantization.} Daily granularity is returned only for
sufficiently short windows. A multi-year daily history must be requested as a
sequence of overlapping chunks and reassembled. The conventional reassembly is
sequential: chunk $2$ is rescaled onto chunk $1$, chunk $3$ onto the rescaled
chunk $2$, and so on. Scale errors then compound multiplicatively along the
chain, and the final segment inherits the accumulated error of every preceding
join.

\paragraph{Non-determinism.} Repeating an identical query on a different day, or
from a different browsing environment, can return different values.
\citet{EichenauerIMS2022} document this instability directly and propose
averaging across repeated samples. Averaging is a substantial improvement over a
single download, but it raises a further question that the applied literature has
largely left open: repeated retrievals are \emph{not} independent samples, so
what exactly does their average estimate, and how much information do $N$
retrievals actually contain?

\subsection{What this paper does}

We treat data acquisition itself as a designed, auditable measurement
experiment rather than as a preliminary inconvenience. The central move is to
refuse the question ``what is the true search volume?'', which the GT system does
not permit us to answer, and to replace it with two questions that are
answerable under stated assumptions:

\begin{enumerate}
\item Under a fixed query specification and a fixed, preregistered retrieval
design, what is the expected value that GT returns for each historical date?
\item On a common scale, what is the latent relative search signal that the
repeated retrievals have in common?
\end{enumerate}

The first estimand is design-based and requires no model of Google's internal
sampling; it is unbiased by construction provided the inclusion rule does not
depend on the realized values. The second is model-based and requires the
assumptions we make explicit in \cref{sec:identification}. Reporting both, and
keeping them distinct, is a deliberate methodological choice: the design-based
quantity is a transparent benchmark against which the model-based reconstruction
can be judged.

\subsection{Contributions}

\begin{enumerate}
\item \textbf{A preregistered acquisition design.} Collection day, collection
stream, technical replicate, chunk order, and time slot are balanced and locked
before any data are seen, and the entire configuration is committed to a
SHA-256 protocol hash (\cref{sec:design}). This makes retrieval-side
researcher degrees of freedom auditable rather than invisible.

\item \textbf{Global overlap-graph calibration.} We replace sequential stitching
with a single weighted least-squares problem on the graph whose nodes are chunks
and whose edges are robustly estimated log-scale differences
(\cref{sec:calibration}). This uses all overlap information at once, removes
error accumulation, and yields identification, connectivity, conditioning, and
cycle-consistency diagnostics that sequential stitching cannot provide.

\item \textbf{Dependence accounting rather than deduplication.} Exact duplicates
are retained in the audit archive and collapsed analytically; near-duplicates are
identified from residuals after the common signal is removed, so that highly
dependent retrievals are not counted as independent evidence
(\cref{sec:duplicates}).

\item \textbf{A three-facet generalizability decomposition.} We estimate variance
components for historical date, collection day, and collection stream and their
interactions, separately for the level, the detection indicator, and the daily
innovation (\cref{sec:gstudy}). The associated D-study answers a question
practitioners actually face: whether to buy reliability with more collection
days, more streams, or more technical replicates.

\item \textbf{A covariance-adjusted consensus with an honest effective sample
size.} Weights solve a minimum-variance problem under the estimated residual
covariance, and we report spectral and Kish effective pull counts so that
dependence is visible in the reported precision (\cref{sec:consensus}).

\item \textbf{Propagation into downstream inference.} The measurement
uncertainty is carried into reliability-corrected OLS and SIMEX rather than
discarded at the boundary of the measurement exercise (\cref{sec:eiv}).
\end{enumerate}

\subsection{Scope and what is deliberately not claimed}

\RACER{} estimates a latent \emph{relative} search signal on a common scale. It
does not estimate absolute Google search counts, and no result in this paper
should be read as identifying them. All unbiasedness, consistency, and
minimum-variance statements are conditional on the assumptions listed in
\cref{sec:identification}, and we state in \cref{sec:scope} exactly which
guarantees survive which assumption failures.

% =============================================================================
\section{The measurement problem}\label{sec:problem}

\subsection{Notation}

Let $t \in \facetT = \{1,\dots,T\}$ index historical calendar dates, and let a
\emph{complete pull} $i$ denote one full reconstruction of the historical
window under the locked query specification. A complete pull consists of $J$
fixed, overlapping chunks; chunk $j$ covers a contiguous window
$W_j \subset \facetT$, and consecutive chunks overlap in at least
$m_{\min}$ days.

Write $Y_{ijt}$ for the raw GT value returned for date $t$ in chunk $j$ of pull
$i$. Each pull is characterized by a design cell: a collection day
$d \in \facetD$, a collection stream $s \in \facetS$, and a technical replicate
$r \in \facetR$.

\begin{remark}[Streams are composite environments]
A stream is a fixed combination of device, browser profile, cookie state,
account state, network route, and IP address. An estimated stream effect is
therefore \emph{not} an IP effect. Attributing it to the IP address alone would
require a crossover factorial design that varies IP while holding the remaining
factors fixed. We are explicit about this because the applied literature
sometimes treats distinct IP addresses as independent sampling units, which the
design does not license.
\end{remark}

\begin{remark}[Why the collection day is the load-bearing facet]\label{rem:cache}
The design treats the collection day as a facet in its own right because the
platform appears to resample on a daily cycle rather than per request.
\citet{Djorno2026} report that repeated identical requests are served from a
cache for twenty-four hours and that a fresh sample is drawn when the cache
resets at midnight UTC.

Two consequences follow. First, repeated retrieval within a day cannot buy
independent information about the sampling error, so the effective rate at which
evidence accumulates is bounded at one fresh draw per day. Any procedure that
reduces sampling variability by averaging repeated pulls --- including the one
recommended by \citet{EichenauerIMS2022} --- therefore trades precision against
calendar time, which is what makes it awkward for real-time work.

Second, the mechanism yields a prediction this framework can be judged by. If the
daily cache holds exactly, the technical-replicate variance component in
\cref{sec:gstudy} should be indistinguishable from zero, and the exact
duplicates retained in \cref{sec:duplicates} are not a nuisance to be cleaned
away but the direct observable signature of the cache. If that component is
bounded away from zero, the daily-cache description is incomplete for the
queries at hand. We state the prediction because it is cheap to test and because
both outcomes are informative; we do not assume it, and no result in this paper
depends on it holding.
\end{remark}

\subsection{Three failure modes of naive construction}

Let $X_t$ denote the latent relative search signal we wish to recover.

\paragraph{(F1) Compounding scale error.} Under sequential stitching with
estimated log-scale corrections $\hat{\ell}_j$, the calibrated value of chunk $J$
carries the error $\sum_{j=2}^{J}(\hat{\ell}_j - \ell_j)$. Even if each join is
unbiased, the variance of the accumulated error grows with the number of joins,
so the oldest or newest part of a long series is systematically less reliable
than the middle.

\paragraph{(F2) Dependence masquerading as replication.} If $N$ retrievals share
a common retrieval-side disturbance, the variance of their mean does not fall as
$1/N$. Reporting $N$ as a sample size overstates precision. In the extreme case
of exact duplicates, $N$ retrievals contain the information of one.

\paragraph{(F3) Attenuation downstream.} If the constructed index enters a
regression as a covariate measured with classical additive error of variance
$\sigma^2_u$, the OLS slope is attenuated by the reliability ratio
$\lambda = \sigma^2_W/(\sigma^2_W + \sigma^2_u)$. Measurement work that stops at
the construction stage leaves this bias in place.

\subsection{What the existing literature establishes}\label{sec:related}

Each failure mode above has been documented separately, and the contribution
here is to place them in one measurement chain rather than to discover them.
\citet{EichenauerIMS2022} identify sampling noise and frequency inconsistency in
GT and propose repeated draws as the remedy; \citet{CebrianDomenech2024} show
that the error in repeated pulls scales with a term's popularity;
\citet{MedeirosPires2021} demonstrate that a single realization can flip a
forecasting conclusion by chance. On construction, \citet{BleherDimpfl2022} knit
multi-annual high-frequency series from overlapping windows, and
\citet{West2020} calibrates across queries with an anchor bank that also handles
rounding and censoring at low volume. \citet{Holzl2025} survey how GT is
misused in the social sciences, including the practice of reading a constructed
index as if it were a count.

Two strands describe the retrieval mechanism itself. \citet{Neumann2023}, working
with Google Health Trends, document the sampling mechanism and recommend
aggregating related phrases with the Boolean \texttt{OR} operator to lift series
above the privacy threshold. Most directly related to the present paper,
\citet{Djorno2026} treat missing values, sampling variability, noise, and trend
as a single preprocessing problem, and report that repeated requests are served
from a cache that resets at midnight UTC, so that a fresh sample is drawn only
once per day. Their remedy is hierarchical clustering, smoothing splines, and
detrending, validated by the accuracy of downstream influenza forecasts.

On empirical comparison this paper is deliberately narrow. Only sequential
stitching is compared directly (\cref{sec:calexp}), and the result is that the
mechanism holds while the accuracy difference does not reach significance.
\citet{West2020}'s anchor bank needs several queries to operate and
\citet{Djorno2026} evaluate against downstream forecast accuracy, neither of
which the single-series, latent-truth simulation here can accommodate. We
therefore make \textbf{no claim} of accuracy superiority over those methods. What
is established is a large improvement over downloading once, and a set of
diagnostics the alternatives do not provide.

What that work leaves open is the step from a smoothed curve to an input for
inference. None of it fixes the retrieval design in advance as a crossed
experiment, decomposes error across collection day and stream, downweights
dependent pulls by an estimated residual covariance, or carries the resulting
uncertainty into the downstream regression. The distinction from
\citet{Djorno2026} is the sharpest available statement of what this paper adds:
that work seeks a cleaner curve and is evaluated by whether forecasts improve,
whereas the target here is the measurement uncertainty itself, which is
estimated, reported, and then propagated. A denoised series with no attached
uncertainty is still consumed downstream as if it were measured without error.
Once it is not, the downstream problem is the generated-regressor problem of
\citet{Pagan1984}: a covariate that is itself an estimate makes the conventional
standard error wrong even when the point estimate is serviceable. \cref{sec:eiv}
therefore treats the constructed series as generated rather than observed, in
the errors-in-variables tradition of \citet{Carroll2006} and
\citet{CookStefanski1994}. Those are the gaps this paper addresses, and
\cref{sec:scope} states precisely which of them come with proofs and which come
with diagnostics only.

\begin{remark}[Why the website version still needs chunking]
Google documents that Trends normalizes a sample of search activity by time and
region and scales it to $0$--$100$, with low volumes reported as
zero~\citep{GoogleFAQ}. The 2025 API alpha offers a consistent scale across
requests, but its historical range is a rolling window of roughly 1{,}800 days
and its values remain relative rather than absolute~\citep{GoogleAPI2025}. A
study spanning 2010--2026 therefore still reconstructs its daily history from
multiple shorter windows, which is what makes the calibration problem in
\cref{sec:calibration} unavoidable rather than optional.
\end{remark}

% =============================================================================
\section{Estimands and identification}\label{sec:identification}

\subsection{Two estimands}

\begin{definition}[Retrieval-mixture expectation]\label{def:mixture}
Fix the query specification $Q$ and a retrieval design that places known
probability $\pi_h$ on design cell $h \in \mathcal{H}$. The
\emph{retrieval-mixture expectation} is
\begin{equation}
\theta_t^{\pi} \;=\; \sum_{h \in \mathcal{H}} \pi_h \,
\E\!\left[Y_{t} \mid Q,\, h \right],
\label{eq:mixture}
\end{equation}
where the inner expectation is taken over the GT system's own randomness
conditional on the cell.
\end{definition}

\begin{definition}[Latent common-scale signal]\label{def:latent}
Let $Z_{it}$ denote the calibrated value of pull $i$ at date $t$ after
within-pull scale alignment and baseline normalization. The \emph{latent
common-scale signal} $X_t$ is defined by the measurement model
\begin{equation}
\Zvec_t \;=\; \one X_t + \evec_t, \qquad \E[\evec_t] = 0, \qquad
\Var(\evec_t) = \Sig .
\label{eq:consensus-model}
\end{equation}
\end{definition}

\Cref{def:mixture} requires no model of the retrieval error and is estimable by
stratified averaging. \Cref{def:latent} requires the assumptions below but
supports variance-optimal estimation.

\begin{remark}[The bridge from $\theta_t$ to $X_t$]\label{rem:bridge}
\Cref{def:mixture} and \cref{def:latent} are not the same target, and two
conditions connect them. First, proportionality: $\theta_t = cX_t$ for some
$c>0$ that does not depend on $t$, so that the expected GT response under a
locked protocol tracks the latent relative signal. This rules out a
systematic algorithmic bias that varies over time, and repeated measurement
cannot test it. Second, negligible calibration error:
$Z_{it}=Y_{it}\exp(-\hat{\ell}_{j(i)})$ is a nonlinear transform of an
\emph{estimate}, and $\E[Y\exp(-\hat{\ell})] \neq a^{-1}\E[Y]$ in general
because $\exp(\cdot)$ is convex. The first-order correction in
\cref{sec:calibration} leaves
$\E[\exp(-\hat{\ell}_j)] = a_j^{-1}\{1+O(n^{-1})\}$ with $n$ the number of
usable overlap days, so $\E[\evec_t]=0$ in \eqref{eq:consensus-model} is an
$O(n^{-1})$ approximation rather than an exact statement. The unbiasedness and
minimum-variance results below should be read under both conditions: the first
needs external validation, the second improves as the overlap design is
densified.
\end{remark}

\begin{remark}[Why no additive pull bias is removed by default]\label{rem:centring}
Version 1.0.0 subtracted $b_i = T^{-1}\sum_t (Z_{it}-m_t)$ from each pull, with
$m_t$ the daily median. Baseline normalization has already set
$T^{-1}\sum_t Z_{it} = 100$ for every pull, so when the baseline spans the full
history --- the default that \texttt{QuerySpec} back-fills --- that centring term
reduces to $b_i = 100 - T^{-1}\sum_t m_t$, the \emph{same constant} for every
$i$. It removes no pull-specific difference; it shifts the entire consensus by
the gap between the mean and the median, and it contradicts $\E[\evec_t]=0$ in
\eqref{eq:consensus-model}. In the controlled simulation it moved mean bias from
machine precision to $-0.1464$ and degraded RMSE in all twenty replications.
Version 1.1.0 therefore defaults \texttt{center\_pull\_bias} to off and reports
\texttt{pull\_bias\_cross\_pull\_sd}, which is near zero exactly when the term
carries no pull-specific information. Enable it only with baseline rescaling
switched off and a genuine additive offset expected.
\end{remark}

\subsection{Assumptions}

\begin{assumption}[Protocol invariance]\label{as:protocol}
Every retrieval in a batch uses the identical query specification, historical
endpoints, and chunk manifest, and the configuration hash is unchanged across
the batch.
\end{assumption}

\begin{assumption}[Value-independent inclusion]\label{as:inclusion}
Whether a retrieval is included in the analysis does not depend on the values it
returned. Retrievals may be excluded for documented operational reasons recorded
before inspection of the values.
\end{assumption}

\begin{assumption}[Positive-overlap multiplicative model]\label{as:overlap}
For a date $t$ present in chunks $j$ and $k$ of the same pull, with both raw
values at or above a prespecified threshold $c_{\min}$,
\begin{equation}
Y_{jt} = a_j X_t \exp(\epsilon_{jt}), \qquad a_j > 0,
\label{eq:chunk-model}
\end{equation}
so that with $\ell_j = \log a_j$ and $\nu_{jkt} = \epsilon_{jt} - \epsilon_{kt}$,
\begin{equation}
r_{jkt} \;=\; \log Y_{jt} - \log Y_{kt} \;=\; \ell_j - \ell_k + \nu_{jkt}.
\label{eq:edge-model}
\end{equation}
\end{assumption}

\begin{assumption}[Connected overlap graph]\label{as:connected}
The graph $G=(V,E)$ with $V$ the set of chunks and $E$ the set of chunk pairs
having at least $m_{\min}$ usable overlapping days is connected.
\end{assumption}

\begin{assumption}[Conditionally centred edge error]\label{as:edge-mean}
$\E[\nu_{jkt} \mid (j,k) \in E] = 0$ for every edge, and the edge-level
estimator is applied to a set of overlap days selected without reference to the
realized differences.
\end{assumption}

\begin{assumption}[Common signal with centred retrieval error]\label{as:common}
Model \eqref{eq:consensus-model} holds with $\Sig$ finite and positive definite
on the retained pulls.
\end{assumption}

\begin{remark}[Where each assumption can fail]
\Cref{as:overlap} fails near zero, because rounding makes the log ratio explode;
this is why only cells at or above $c_{\min}$ contribute to edges, while zeros
are never imputed to positive values. One consequence of the aggregation rule
must be stated plainly: when a date is covered by several chunks and only some
of them return zero, the weighted average is positive, so that date is no longer
zero in the final series. Zeros survive only where every covering chunk is zero.
That is a corollary of averaging rather than an imputation decision, but it
changes what detection reliability is measured on, so
\texttt{CalibrationResult.\allowbreak diagnostics} reports
\texttt{n\_dates\_with\_zero\_chunk} and
\texttt{n\_dates\_zero\_masked\_by\_aggregation} for audit. \Cref{as:connected} fails if the chunk design leaves a gap; the
implementation raises an error rather than silently returning a partially
identified series. \Cref{as:common} fails if a subset of pulls shares a
retrieval-side bias that is not shared by the rest; this is precisely what the
near-duplicate diagnostics in \cref{sec:duplicates} are designed to expose.
\end{remark}

% =============================================================================
\section{Stage 0: preregistered acquisition design}\label{sec:design}

The design fixes, before collection begins: the query specification; the
historical endpoints; the chunk manifest (window length, step, minimum overlap);
the set of collection-day offsets; the set of streams; the number of technical
replicates; the time slots; and the acceptance thresholds. Stream order is
rotated across collection days and chunk order is rotated across pulls, so that
order effects are balanced rather than confounded with the facets of interest.

The entire configuration is serialized canonically and hashed:
\begin{equation}
h \;=\; \mathrm{SHA256}\big\{\mathrm{canonical}(Q, \facetD, \mathcal{C}, \mathcal{T})\big\},
\label{eq:hash}
\end{equation}
where $\mathcal{C}$ collects the calibration settings and $\mathcal{T}$ the
decision thresholds. Every raw observation may store $h$. Any change to a locked
setting changes $h$, so configuration drift between the design and the analysis
is detectable rather than assumed away.

\paragraph{Pilot separation.} An initial monitoring window may be used to check
for gaps, query drift, exact duplicates, and connectivity. If that inspection
leads to any change in thresholds, stitching rules, query, code, or decision
tree, the monitored data must be relabelled as pilot and the confirmatory batch
restarted. Otherwise the same data are used both to design and to validate the
acceptance rule, producing data-dependent acceptance bias. For the same reason,
all GT processing rules must be locked before any return, volatility, spread, or
VaR result is examined; selecting the construction that gives the most
significant financial result converts measurement validation into
outcome-guided model selection.

% =============================================================================
\section{Stage 1: global overlap-graph calibration}\label{sec:calibration}

\subsection{Robust edge estimation}

For each chunk pair $(j,k)$ with at least $m_{\min}$ usable overlap days, the
observed log ratios $\{r_{jkt}\}$ are summarized by a Huber M-estimator of
location \citep{Huber1964}, which solves
\begin{equation}
\sum_{t} \psi_c\!\left(\frac{r_{jkt} - \mu}{\hat{\sigma}}\right) = 0,
\qquad
\psi_c(u) = \max\{-c, \min(c, u)\},
\label{eq:huber}
\end{equation}
with $c = 1.345$ and $\hat{\sigma}$ a robust scale. The estimator returns the
location $\hat{r}_{jk}$, a robust scale, and a variance $\hat{v}_{jk}$ that
becomes the edge precision. Robustness matters here because a small number of
overlap days can be contaminated by revisions or by the rounding boundary, and a
sample mean of log ratios is not resistant to those points.

\subsection{Graph weighted least squares}

Collect the edges into a design matrix $B$ with rows indexed by edges: row $e$
for edge $(j,k)$ has $+1$ in column $j$, $-1$ in column $k$, and zero elsewhere.
Because only differences are identified, one chunk is fixed as the reference by
setting $\ell_{\mathrm{ref}} = 0$; the implementation selects the
highest-degree node, which minimizes the leverage of any single poorly estimated
edge. With $\bm{r}$ the vector of edge estimates and
$\Omega = \diag(\hat{v}_e)$,
\begin{equation}
\hat{\bm{\ell}} \;=\; \argmin_{\bm{\ell}}\;
(\bm{r} - B\bm{\ell})^{\top}\Omega^{-1}(\bm{r} - B\bm{\ell})
\;=\; (B^{\top}\Omega^{-1}B)^{-1}B^{\top}\Omega^{-1}\bm{r}.
\label{eq:gwls}
\end{equation}

\begin{proposition}[Identification]\label{prop:ident}
Under \cref{as:connected}, $B$ has rank $|V| - 1$ after the reference
normalization, so $\hat{\bm{\ell}}$ in \eqref{eq:gwls} is uniquely defined.
\end{proposition}

\begin{proof}
$B$ is the incidence matrix of $G$. For a graph with $|V|$ nodes, the rank of
the incidence matrix equals $|V|$ minus the number of connected components.
Connectivity gives one component, hence rank $|V|-1$. Deleting the reference
column removes exactly the one-dimensional null space spanned by $\one$, leaving
a full-column-rank matrix and a positive definite normal matrix
$B^{\top}\Omega^{-1}B$.
\end{proof}

\begin{proposition}[Conditional unbiasedness and efficiency]\label{prop:gwls}
Under \cref{as:overlap,as:connected,as:edge-mean}, and treating $\Omega$ as
known, $\hat{\bm{\ell}}$ is unbiased for $\bm{\ell}$ and is the minimum-variance
estimator within the class of linear unbiased estimators of $\bm{\ell}$ based on
the edge estimates.
\end{proposition}

\begin{proof}
Write $\bm{r} = B\bm{\ell} + \bm{\nu}$ with $\E[\bm{\nu}] = 0$ by
\cref{as:edge-mean} and $\Var(\bm{\nu}) = \Omega$. This is a linear model with
full-column-rank design after normalization, so the Gauss--Markov theorem
applies directly to the weighted least-squares estimator \eqref{eq:gwls}.
\end{proof}

That proposition says only that graph WLS is best among linear estimators built
on the edge estimates. It does not answer the question the design actually turns
on --- \emph{how much better than sequential stitching} --- and failure mode F1
was until now an argument in words. The next proposition makes it an identity,
and one that depends on the chunk design alone rather than on how noisy the
retrievals turn out to be.

\begin{proposition}[Calibration variance is an effective resistance]\label{prop:reff}
Read the overlap graph as an electrical network in which edge $e$ carries a
\emph{conductance} equal to its precision, $w_e = 1/\Var(\hat\delta_e)$. Then for
any chunk $j$ and reference $r$,
\begin{equation}
\Var(\hat\ell_j - \hat\ell_r \mid B) = R_{\mathrm{eff}}(j, r),
\label{eq:reff}
\end{equation}
the effective resistance between $j$ and $r$ under the weighted Laplacian
$L = B'WB$. Two consequences follow.

\textbf{(i)} Sequential stitching is the same estimator restricted to a spanning
path, whose effective resistance is the plain \emph{sum} of the resistances it
walks and therefore grows with the number of joins. That is F1, stated exactly
rather than by analogy.

\textbf{(ii)} (Rayleigh monotonicity) Adding an edge cannot raise the effective
resistance between any pair of nodes. Hence
\begin{equation}
R_{\mathrm{eff}}^{\text{graph}}(j, r) \;\le\; R_{\mathrm{eff}}^{\text{path}}(j, r)
\qquad \text{for every } j,
\label{eq:rayleigh}
\end{equation}
so \textbf{global graph calibration is never worse than any sequential stitch in
calibration variance}, and the margin can be computed from the design
\emph{before} any data is collected.
\end{proposition}

\begin{proof}
Grounding at $r$ and deleting its row and column leaves the grounded Laplacian
$B_{-r}'WB_{-r}$, whose inverse is the estimator covariance in \eqref{eq:gwls}.
In network theory the $j$th diagonal entry of a grounded Laplacian inverse is the
effective resistance from $j$ to ground, which is \eqref{eq:reff}. Consequence
(i) is resistances in series. Consequence (ii) is Rayleigh's monotonicity
theorem: an added edge is an added parallel path, which cannot increase the
resistance between any two points.
\end{proof}

\begin{remark}[F1 becomes a quantity the design stage can compute]\label{rem:designtool}
Since $\Var(\hat\delta_e) \approx \sigma^2/n_e$ for an overlap of $n_e$ days with
per-day log-ratio dispersion $\sigma$, the $\sigma^2$ cancels in the ratio of
path to graph: \textbf{the variance reduction is a function of the chunk design
alone}. Converting to a terminal scale error at the $\sigma = 0.0188$ measured on
real data:

\begin{table}[H]
\centering
\small
\caption{Calibration error implied by a chunk design through \eqref{eq:reff},
under $\Var(\hat\delta_e)=\sigma^2/n_e$ at the median dispersion measured on one
year of one keyword. \textbf{Only the first row can be checked against data; the
rest are extrapolations} --- see \cref{rem:extrapolation}.}
\label{tab:designreff}
\begin{tabular}{lrrrr}
\toprule
Design (span, window/step) & Chunks & Reduction & Sequential & Graph\\
\midrule
\multicolumn{5}{l}{\emph{Measured, from the estimated edge variances}}\\
1 year, 180/30 & 8 & $\mathbf{9.13\times}$ & --- & ---\\
\midrule
\multicolumn{5}{l}{\emph{Extrapolated, from the idealized edge model}}\\
1 year, 180/30 & 8 & $10.5\times$ & $0.40\%$ & $0.12\%$\\
5 years, 180/30 & 56 & $18.2\times$ & $1.15\%$ & $0.27\%$\\
15 years, 180/30 & 178 & $20.1\times$ & $2.06\%$ & $0.46\%$\\
15 years, 180/15 (recommended) & 355 & $140.5\times$ & $2.79\%$ & $0.23\%$\\
\bottomrule
\end{tabular}
\end{table}

\begin{remark}[What in this table is measured and what is not]\label{rem:extrapolation}
The identity \eqref{eq:reff} and the ordering \eqref{eq:rayleigh} are theorems and
need no data. The \emph{numbers} in the table are a different matter, and two of
their assumptions are contradicted by the same year of data they are calibrated
on. Both contradictions run in the direction that flatters the conclusion, so
they are stated rather than left for a reader to find.

\textbf{The idealized edge model overstates the reduction by about 15\%.} On the
one design where both can be computed, the estimated edge variances give
$9.13\times$ against the model's $10.5\times$: five of the 26 overlaps are
degenerate and take the variance floor rather than $\sigma^2/n_e$, and the Huber
weights leave $n_{\mathrm{eff}}<n_e$. The extrapolated rows have no such check.

\textbf{$\sigma$ is not a constant.} Across the 21 informative overlaps it ranges
over a factor of $1.75$ \emph{within a single year} and correlates $-0.760$
($p=0.0001$) with the value level of the pair, which is what the rounding floor
$1/(\sqrt{12}\,v)$ predicts. A fifteen-year series spans far wider variation in
search volume than one year of it does, so the low-volume eras carry a larger
$\sigma$ and the absolute percentages in the last three rows are understated. The
\emph{ratio} is more robust, because $\sigma$ largely cancels between path and
graph, but it does not cancel exactly once $\sigma$ varies across edges.

The reach is also worth stating plainly: one year, one keyword, one pull,
extrapolated fifteen-fold. \textbf{The claim that F1 is material at long horizons
is a prediction of \cref{prop:reff}, not a measurement.} Testing it needs a
long-horizon collection, which \cref{sec:design} specifies and this
manuscript has not carried out.
\end{remark}

From version 1.4.0 \texttt{CalibrationResult.diagnostics} reports
\texttt{max\_\allowbreak log\_\allowbreak scale\_\allowbreak variance},
\texttt{sequential\_\allowbreak max\_\allowbreak log\_\allowbreak scale\_\allowbreak variance},
and their ratio
\texttt{calibration\_\allowbreak variance\_\allowbreak reduction}.
\end{remark}

\begin{remark}[Why this dominates sequential stitching]
Sequential stitching is the special case in which only the $J-1$ consecutive
edges are used and each is applied in turn. \Cref{prop:gwls} uses every edge,
including the non-consecutive overlaps that a long window and a short step
generate, and weights each by its estimated precision. It also makes error
accumulation impossible: no estimate is defined as a function of a previously
calibrated estimate. In addition, the residual $\bm{r} - B\hat{\bm{\ell}}$ has a
direct interpretation as cycle inconsistency, which is reported as a diagnostic
and has no analogue in the sequential procedure.
\end{remark}

\subsection{Reconstruction and normalization}

Calibrated observations are $\tilde{Y}_{jt} = Y_{jt}\exp(-\hat{\ell}_j)$,
optionally with a lognormal bias correction. Overlapping dates are aggregated by
inverse-variance weighting. The resulting series is normalized so that its mean
over the prespecified baseline window equals $100$. This normalization defines
the scale of the estimand; it is not an estimate of absolute search volume.

% =============================================================================
\subsection{A direct comparison with sequential stitching}\label{sec:calexp}

Failure mode F1 argued that sequential stitching accumulates scale error along
the chain and that solving the graph avoids it. Until now that was an argument,
not a measurement. \texttt{src/racergt/baselines.py} implements sequential
stitching in the style of Bleher and Dimpfl's knitting, reusing the \emph{same}
Huber edge estimator, the same $c_{\min}$ gate, and the same baseline
normalization; the only difference is that it walks a spanning path instead of
solving the whole graph, so any gap is attributable to that choice alone.
\texttt{examples/run\_calibration\_comparison.py} runs both.

\begin{table}[H]
\centering
\small
\caption{Global graph calibration against sequential stitching, 20 seeds by 2
pulls per cell. Advantage is sequential minus graph in RMSE, so positive favours
the graph.}
\label{tab:calexp}
\begin{tabular}{lrrrrrrl}
\toprule
Design & Joins & Graph & Sequential & Advantage & MCSE & $p$ & Graph wins\\
\midrule
60-day step & 5 & 7.4827 & 7.4677 & $-0.0150$ & 0.0053 & $0.008$ & 13/40\\
30-day step & 9 & 7.3026 & 7.2967 & $-0.0059$ & 0.0049 & $0.234$ & 15/40\\
15-day step & 17 & 7.2085 & 7.2125 & $+0.0040$ & 0.0072 & $0.576$ & 15/40\\
15-day step, low noise & 17 & 1.9245 & 1.9356 & $+0.0111$ & 0.0066 & $0.098$ & 21/40\\
\bottomrule
\end{tabular}
\end{table}

\textbf{The mechanism holds; its effect is second order.} The accumulation F1
describes is real and its magnitude is clear: under the 15-day design the
cross-replication standard deviation of the recovered log scale rises from
$0.0259$ early in the chain to $0.1228$ at the end, a factor of $4.75$, while
the graph's ratio is $1.12$. Error does accumulate along a spanning path, and
solving jointly does prevent it.

It does not translate into a detectable advantage in reconstruction error. The
direction matches F1 --- going from 5 joins to 17 moves the advantage from
$-0.0150$ to $+0.0040$, and cutting retrieval noise from $0.06$ to $0.01$, so
that calibration error carries a larger share of the total, raises it further to
$+0.0111$ --- but even in that setting, purpose-built to favour F1, $p=0.098$ with
21 wins out of 40. At the shortest chain the graph is significantly \emph{worse}
($p=0.008$), because its reference chunk is chosen by highest degree rather than
fixed at the first chunk, and because it additionally applies the lognormal
correction.

\begin{remark}[Three mechanisms, one conclusion]\label{rem:pattern}
This is the third result of the same shape. Covariance-adjusted weighting
correctly identifies and down-weights dependent pulls (weights correlate
$-0.562$ with shared residual dependence) without improving RMSE. The prediction
that heterogeneous dependence would change that was tested and failed. Global
graph calibration does prevent error accumulating along the chain ($1.12$ against
$4.75$) and likewise does not improve RMSE.

The common structure is that each mechanism works as designed and can be shown
to work, while its effect on \emph{point accuracy} is second order and swamped
by retrieval noise. For a user that conclusion matters more than any single
number: the value of this framework is not that it produces a more accurate
series, but that it produces a series whose quality can be judged. Effective
pull counts, dependence structure, cycle-consistency statistics, per-day
conditional standard errors, and propagatable measurement uncertainty are
information that RMSE neither expresses nor replaces.
\end{remark}

% =============================================================================
\subsection{Calibration on real Google Trends data}\label{sec:realdata}

All of the evidence above comes from the package's own simulator, which draws
from the same proportionality model the estimator assumes. That is circular: it
can show that the estimator solves the problem it was handed, never that the
problem is the right one. This section runs the calibration stage on real
downloads --- keyword \texttt{bitcoin}, United States, daily, calendar year 2024,
eight chunks on 180-day windows with a 30-day step, 1{,}440 observations. The
converter is \texttt{examples/\allowbreak ingest\_\allowbreak trends\_\allowbreak csv.py} and the analysis is
\texttt{examples/\allowbreak run\_\allowbreak real\_\allowbreak data\_\allowbreak calibration.py}.

\textbf{What this data cannot test, stated first.} It is a \emph{single pull}:
one collection day, one stream, one replicate. Every cross-pull mechanism is
therefore \emph{entirely} untested here --- covariance-adjusted consensus
weighting, exact and near-duplicate diagnosis, the three-facet variance
components, effective pull counts, and the RMSE comparison of
\cref{sec:calexp} all require repeated collection across days and environments.
This section does not touch them and claims nothing about them. Testing the full
design requires the month-long collection schedule of \cref{sec:design}. Real
data also carries no latent truth, so \textbf{reconstruction error cannot be
computed} and neither method can be called more accurate.

\subsubsection{The normalization creates zero-information overlaps}

The first finding is one the simulator cannot structurally produce. Trends
divides each chunk by the maximum inside \emph{that request's own window}. When
two windows contain the same peak day they are divided by the same number and,
after rounding, are \emph{identical cell for cell} wherever they overlap. Here
C0001--C0003 all peak on 2024-03-05, C0004--C0005 on 2024-08-05, and
C0007--C0008 on 2024-12-05, so 5 of the 26 overlap edges have a log ratio that
is identically zero with no dispersion: the eight chunks form only \textbf{four
normalization groups}.

Such an edge carries no scale information, yet it still enters the weighted
least-squares fit at whatever weight \texttt{edge\_\allowbreak variance\_\allowbreak floor} and
\texttt{max\_\allowbreak edge\_\allowbreak weight} permit --- numerical safety constants, not a
statistical argument. \texttt{simulate\_\allowbreak racergt\_\allowbreak data} multiplies every chunk by
continuous lognormal noise \emph{before} taking the window maximum, and at its
default settings the chance that two chunks round to identical integers
throughout an overlap is negligible: zero degenerate edges in 520 simulated,
against 5 in 26 real ones. (Setting \texttt{chunk\_\allowbreak noise\_\allowbreak sd} to $0$ would let the
simulator produce them, but that is neither its default nor the setting under
which the earlier evidence was generated.) Version 1.4.0 therefore adds
\texttt{n\_\allowbreak zero\_\allowbreak dispersion\_\allowbreak edges}, \texttt{n\_\allowbreak informative\_\allowbreak edges} and
\texttt{n\_scale\_groups} to \texttt{CalibrationResult.\allowbreak diagnostics}. They
\textbf{report only and change no estimate}; what to do about a degenerate edge
is the analyst's judgment, not the package's.

The rate depends on the series' peak structure and \textbf{cannot be
extrapolated from one keyword}. The mechanism is deductive --- a shared window
maximum implies an identical overlap --- but the proportion is an empirical
question.

\subsubsection{The proportionality assumption survives}

The whole calibration stage rests on overlapping chunks being proportional up to
a constant. The overlaps run to 150 days, so the assumption is heavily
over-identified and can be checked directly.

\begin{table}[H]
\centering
\small
\caption{Calibration diagnostics on real data. Bitcoin, United States, daily,
2024; eight chunks, one pull.}
\label{tab:realdata}
\begin{tabular}{lr}
\toprule
Diagnostic & Value\\
\midrule
Overlap graph connected & yes\\
Edges (a spanning path uses 7) & 26\\
Independent cycles & 19\\
Testable triangles & 45\\
Normal matrix condition number & 10.82\\
Weighted edge RMSE & 0.001379\\
\midrule
Zero-dispersion edges & 5 / 26\\
Normalization groups & 4 (among 8 chunks)\\
\midrule
Median triangle closure error & 0.001783 log points\\
Maximum triangle closure error & 0.011769 log points (1.184\% scale error)\\
Overlaps with drift $p<0.05$ & 1 / 21 (1.05 expected by chance)\\
Largest residual mis-scaling & $-0.032\%$ \emph{(no power; see \cref{rem:power})}\\
\bottomrule
\end{tabular}
\end{table}

Triangle closure is one of the framework's central advantages over sequential
stitching, and \textbf{stitching cannot supply it at all}: the model forces
$d_{ij}+d_{jk}=d_{ik}$ for any three mutually overlapping chunks, while a
spanning path uses only the edges it walks and can neither see nor report a
violation. Across 45 triangles the median closure error is $0.0018$ log points
and the largest is $0.0118$, a $1.18\%$ scale error. Only 1 of 21 informative
overlaps shows drift at $p<0.05$, against $1.05$ expected by chance.

\begin{remark}[What these diagnostics can and cannot detect]\label{rem:power}
A diagnostic that passes is worth nothing until it is shown capable of failing.
Injecting a known perturbation into one chunk and re-running gives each test a
minimum detectable effect.

\begin{table}[H]
\centering
\small
\caption{Response of each diagnostic to a known perturbation of chunk C0005.
Baseline maximum closure error is $0.0118$; baseline drift count is 1 of 21.}
\label{tab:power}
\begin{tabular}{lccc}
\toprule
Perturbation & Max closure error & Drift $p<0.05$ & Per-chunk deviation\\
\midrule
Uniform rescale, any size & $0.0118$ (none) & 1 (none) & none\\
Trend, $2\%$ per 100 days & $0.0232$ & 7 of 22 & none\\
Nonlinear, $2\%$ & $0.0118$ (none) & 5 of 22 & none\\
Nonlinear, $5\%$ & $0.0212$ & 5 of 22 & none\\
\bottomrule
\end{tabular}
\end{table}

\textbf{The per-chunk deviation test has exactly zero power and is reported only
to say so.} Multiplying a chunk by a constant is absorbed entirely into
$\hat\ell_j$, leaving $Y_{jt}\exp(-\hat\ell_j)$ unchanged, so injected errors of
$2\%$ and $20\%$ return the same number to the digit. That is not a weakness of
the estimator --- a uniform rescale is the parameter, not a violation --- but it
means the $-0.032\%$ in \cref{tab:realdata} is evidence of nothing.

Cycle closure and the drift test do have power, and against violations that a
scale cannot absorb. The drift test is the more sensitive of the two, catching a
$2\%$ nonlinear distortion that closure misses. Reading the baseline against the
table: the observed maximum closure error is half what a $2\%$-per-100-days trend
would produce, so \textbf{violations of that size or larger are excluded}, which
is the strongest statement this data supports. It is not the same as the model
being correct.
\end{remark}

\subsubsection{The simulator overstates chunk noise by about 4.5 times}

This data cannot test F1, but it decisively tests something else: whether the
simulator describes reality. \texttt{SimulationSettings.\allowbreak chunk\_\allowbreak noise\_\allowbreak sd}
defaults to $0.03$, multiplying every chunk on every day by an independent
lognormal disturbance before the maximum is taken. The real overlaps have a
median log-ratio dispersion of $0.0188$.

Subtracting a closed-form integer-rounding floor returns $0$, but that floor
assumes the two chunks round \emph{independently}. They do not: they round the
same numbers differing only by a scale factor, so their rounding errors are
correlated, the closed form overstates itself, and the recovered noise is biased
towards zero --- the direction that flatters this conclusion. It is therefore not
used. A parametric bootstrap replaces the assumption with the mechanism: take the
calibrated daily series, multiply by chunk noise at a stated level, divide by the
window maximum, round, and measure dispersion exactly as on the real chunks.

\begin{table}[H]
\centering
\small
\caption{Sweeping \texttt{chunk\_\allowbreak noise\_\allowbreak sd} against the observed overlap
dispersion of $0.018814$. Two hundred replications per level.}
\label{tab:chunknoise}
\begin{tabular}{lrr}
\toprule
$\sigma_{\text{chunk}}$ & Median overlap dispersion & Ratio to observed\\
\midrule
$0$ (rounding only) & 0.017757 & 0.94\\
$0.005$ & 0.017644 & 0.94\\
$0.0066$ (\textbf{matched}) & 0.018814 & 1.00\\
$0.010$ & 0.021892 & 1.16\\
$0.020$ & 0.033624 & 1.79\\
$0.030$ (\textbf{simulator default}) & 0.046424 & 2.47\\
\bottomrule
\end{tabular}
\end{table}

Integer rounding alone accounts for $94\%$ of the observed dispersion, and the
matched value is $\sigma_{\text{chunk}}\approx 0.0066$, about $4.5$ times smaller
than the simulator's default. Below about $0.005$ the curve is flat within Monte
Carlo error --- rounding dominates and the parameter is not identified there ---
but $0.03$ produces $2.47$ times the observed dispersion and is firmly excluded.

The sweep itself is checked the same way the diagnostics are in \cref{rem:power},
by giving it something to find. Injecting a known disturbance into the real
chunks, renormalizing the way Trends does, and rerunning recovers $0.0067$,
$0.0099$, $0.0234$ and $0.0312$ at injected levels of $0$, $0.01$, $0.02$ and
$0.03$. The last is what matters: had the simulator's default been the truth, the
procedure would have returned about $0.03$ rather than the small value it returns
on untouched data, so the gap is not the estimator failing.

This \emph{strengthens} rather than undermines \cref{sec:calexp}. That comparison
was run in a regime where calibration noise was several times too large, and even
in that F1-friendly regime the global graph secured no detectable accuracy
advantage. Rerunning it at realistic noise would only shrink the gap.

\subsubsection{Why one year cannot see F1, and when F1 starts to matter}

The two methods are nearly indistinguishable on real data: correlation
$1.000000$, median daily relative difference $0.00026$, maximum $0.00123$, and a
largest absolute difference of $0.195$ points on the 0--100 index scale.

\begin{remark}[A methodological error worth naming]\label{rem:n1}
Testing F1 by asking whether the gap between the two point estimates grows along
the chain is the \emph{wrong estimator}. F1 is a claim about a variance, and a
single pull supplies one realization; a variance cannot be estimated from
$n = 1$, so that test has no power by construction, at any sample size.
\Cref{prop:reff} removes the need to sample at all: the variance ratio is an
effective-resistance ratio, a deterministic function of the design.
\end{remark}

Applying \cref{prop:reff} to this design, the graph's largest
$\Var(\hat\ell_j)$ is $1.01 \times 10^{-6}$ against $9.26 \times 10^{-6}$ for the
spanning path, a reduction of $\mathbf{9.13}$ times, and the reduction grows
monotonically with position along the chain ($2.1\times$ to $9.1\times$).
\textbf{F1's mechanism is operating}, and its size follows from the design
without any test.

In absolute terms, however, the terminal error is $0.40\%$ for sequential
stitching and $0.12\%$ for the graph. \textbf{A large ratio and a small
magnitude} are not in tension, and together they explain both why one year sees
nothing and why the three earlier simulation experiments returned nulls.

\begin{remark}[F1 is predicted to be material at the recommended protocol]\label{rem:f1scale}
\Cref{tab:designreff} extrapolates the error out along the historical span. At 15
years with 180-day windows and a 15-day step --- the specification recommended
here --- sequential stitching would carry a terminal scale error near $2.79\%$
against $0.23\%$ for the graph. On an attention index spanning fifteen years, a
level error of that size at one end is a slow drift hard to separate from a
trend: exactly the artificial low-frequency drift the Google Trends literature
warns about. \Cref{rem:extrapolation} states what those two numbers rest on; they
are a prediction of \cref{prop:reff} and have not been measured.

The correct statement about F1 is therefore that the mechanism is real and
provable (\cref{prop:reff}), that it is \emph{predicted} to be material at the
recommended scale, and that the one year measured here cannot see it because the
design's absolute error ceiling is $0.40\%$. The nulls in \cref{sec:calexp} and \cref{sec:monte-carlo} are not
evidence against F1; they are what \cref{prop:reff} predicts, since on short
chains the absolute difference between effective resistances must be small.

This revises \cref{rem:pattern}. Its reading was that each mechanism works while
its effect is second order. \Cref{prop:reff} shows that for calibration the
effect is not second order but \emph{scale-dependent} --- it grows with the
number of joins, and every experiment so far ran at the end of the range where it
is smallest. Beyond the diagnostics in \cref{tab:realdata} that a spanning path
structurally cannot supply, the reason to prefer global calibration is the
unconditional guarantee in \eqref{eq:rayleigh}: it is never worse.
\end{remark}

% =============================================================================
\section{Stage 2: duplicate and dependence diagnostics}\label{sec:duplicates}

Two retrievals may be numerically identical, or nearly so. Both cases violate
the implicit assumption that repeated pulls contribute independent information,
but they require different treatment.

\paragraph{Exact duplicates.} Vectors are hashed after rounding to a fixed
number of decimals. Exact groups are retained in the audit archive---deleting
them would destroy evidence about the retrieval process---and are collapsed
analytically before the consensus step, with multiplicity recorded.

\paragraph{Near-duplicates.} Raw correlation between two pulls is uninformative,
because all pulls share the same strong latent signal and are therefore
correlated at close to one by construction. The diagnostic instead removes the
common signal first and examines the residual correlation, together with
value-anchored measures: mean absolute deviation on the $0$--$100$ scale, exact
cell agreement, and the Jaccard index of the positive-value sets. Pairs that
exceed the prespecified thresholds form the edges of a dependence graph.

\begin{remark}[Components are diagnostic, not equivalence classes]
Connected components of the dependence graph are reported as connectivity sets.
Membership in a component does \emph{not} imply that every pair inside it
exceeds the threshold, because connectivity is transitive while the pairwise
rule is not. The acceptance rule uses the share of the batch occupied by the
largest component, which is the quantity that actually matters for whether the
batch is dominated by one dependence cluster.
\end{remark}

% =============================================================================
\section{Stage 3: generalizability decomposition}\label{sec:gstudy}

Classical reliability coefficients assign all non-signal variation to a single
undifferentiated error term. Generalizability Theory
\citep{Cronbach1972,Brennan2001} instead decomposes it across the facets of the
measurement design, which is what a researcher needs in order to decide where to
spend the next unit of collection effort.

For a fully crossed design over historical date $t$, collection day $d$, stream
$s$, and replicate $r$,
\begin{equation}
Y_{tdsr} = \mu + T_t + D_d + S_s + (TD)_{td} + (TS)_{ts} + (DS)_{ds}
+ (TDS)_{tds} + E_{tdsr}.
\label{eq:gstudy}
\end{equation}
Variance components are estimated by the method of moments from the balanced
ANOVA sums of squares. The implementation validates full crossing and equal cell
counts before estimating, and raises an error otherwise, rather than silently
producing components from an unbalanced design.

\begin{remark}[Confounding with a single replicate]
When $n_r = 1$, $E$ cannot be separated from $(TDS)$; the implementation retains
the combined term and reports it as such. This is stated rather than hidden
because it changes the interpretation of the residual component.
\end{remark}

The decomposition is run on three transformations, which answer three different
questions:
\begin{itemize}
\item \textbf{level}, $Y$ itself: do relative magnitudes reproduce?
\item \textbf{detection}, $\indic{Y > 0}$: does the zero/non-zero pattern
reproduce? This is the binding constraint for low-volume terms.
\item \textbf{innovation}, $\Delta\log(1+Y)$: do daily \emph{changes} reproduce?
This is the relevant question when the index enters a model in differences,
and it is systematically harder to satisfy than the level criterion.
\end{itemize}

The generalizability coefficient for relative decisions and the dependability
coefficient for absolute decisions follow in the usual way, and the D-study
evaluates them on a grid of hypothetical designs, so the marginal value of an
additional collection day can be compared with that of an additional stream.

% =============================================================================
\section{Stage 4: the design-based reference estimator}\label{sec:designest}

Let $h$ index the prespecified collection-day $\times$ stream cells with target
weights $\pi_h$ summing to one, and let $n_h$ be the number of completed pulls
observed in cell $h$. The stratified estimator is
\begin{equation}
\hat{\theta}^{\pi}_t \;=\; \sum_{h} \pi_h \,
\frac{1}{n_h}\sum_{i \in h} Z_{it}.
\label{eq:design-est}
\end{equation}

\begin{proposition}[Design unbiasedness]\label{prop:design}
Under \cref{as:protocol,as:inclusion}, and provided $n_h \ge 1$ for every cell
with $\pi_h > 0$, $\hat{\theta}^{\pi}_t$ is unbiased for the retrieval-mixture
expectation $\theta^{\pi}_t$ of \cref{def:mixture}.
\end{proposition}

\begin{proof}
Conditional on cell $h$, the retrievals in that cell are draws from the GT
system under a fixed specification, so
$\E[n_h^{-1}\sum_{i \in h} Z_{it}] = \E[Y_t \mid Q, h]$ by
\cref{as:inclusion}, which guarantees that inclusion carries no information
about the values. Taking the $\pi$-weighted sum over cells and applying
linearity gives \eqref{eq:mixture}.
\end{proof}

\begin{remark}[Dependence affects the variance, not the bias]
\Cref{prop:design} does not require the retrievals to be independent. Dependence
across pulls inflates $\Var(\hat{\theta}^{\pi}_t)$ but leaves the expectation
unchanged. This is exactly why the estimator is retained as a transparent
reference: its validity rests on the design, not on a covariance model that
could itself be misspecified.
\end{remark}

% =============================================================================
\section{Stage 5: covariance-adjusted consensus}\label{sec:consensus}

\subsection{Minimum-variance weights}

Under \cref{as:common}, consider linear estimators
$\hat{X}_t = \wvec^{\top}\Zvec_t$ with $\wvec^{\top}\one = 1$, which is the
condition for unbiasedness.

\begin{proposition}[BLUE weights]\label{prop:blue}
Under \cref{as:common} with $\Sig$ known and positive definite,
\begin{equation}
\wvec^{*} \;=\; \frac{\Sig^{-1}\one}{\one^{\top}\Sig^{-1}\one}
\label{eq:gls-weights}
\end{equation}
solves $\min_{\wvec}\, \wvec^{\top}\Sig\wvec$ subject to
$\wvec^{\top}\one = 1$, and the resulting
$\hat{X}_t = \wvec^{*\top}\Zvec_t$ has variance
$(\one^{\top}\Sig^{-1}\one)^{-1}$, the minimum over all linear unbiased
estimators.
\end{proposition}

\begin{proof}
The Lagrangian is
$L = \wvec^{\top}\Sig\wvec - 2\lambda(\wvec^{\top}\one - 1)$ with first-order
condition $\Sig\wvec = \lambda\one$, hence $\wvec = \lambda\Sig^{-1}\one$.
Imposing $\wvec^{\top}\one=1$ gives
$\lambda = (\one^{\top}\Sig^{-1}\one)^{-1}$ and \eqref{eq:gls-weights}. The
objective is strictly convex because $\Sig \succ 0$, so the stationary point is
the unique minimum. Substituting yields the stated variance.
\end{proof}

\subsection{Stabilization in finite samples}

\Cref{prop:blue} assumes $\Sig$ known. In practice the number of pulls is of the
same order as, or smaller than, the number of dates used to estimate the
covariance, so the sample covariance is ill-conditioned and its inverse is
unstable. The implementation therefore:

\begin{enumerate}
\item collapses exact duplicates analytically, so that a repeated vector cannot
inflate the apparent rank;
\item removes a pull-level offset, so that a constant retrieval bias is not
mistaken for signal disagreement;
\item estimates $\Sig$ by Ledoit--Wolf shrinkage \citep{LedoitWolf2004} toward a
well-conditioned target;
\item imposes $w_i \ge 0$ and $w_i \le \bar{w}$ while retaining
$\sum_i w_i = 1$.
\end{enumerate}

\begin{remark}[What the constraints cost]
The nonnegativity and cap constraints are stabilization devices, not
improvements on \cref{prop:blue}. The strict BLUE property belongs to the
unconstrained solution. The constrained solution is a minimum-variance convex
combination: it gives up the theoretical optimum in exchange for refusing the
large offsetting positive and negative weights that an ill-conditioned
$\hat{\Sig}^{-1}$ would otherwise produce. We regard this as the right trade for
a measurement instrument that must behave predictably on data the analyst has
not yet seen.
\end{remark}

\subsection{Effective number of pulls}

Because dependence is the central concern, the raw count of pulls is not
reported as a sample size. We report instead the Kish effective count
\citep{Kish1965} $1/\sum_i w_i^2$ and a spectral effective count derived from the
eigenvalues of the residual correlation matrix. Both fall below the nominal
count when retrievals are dependent, and the acceptance rule in
\cref{sec:decision} is stated in terms of the spectral count.

The spectral count carries a known downward bias that should be read alongside
it. Residuals are formed by subtracting the daily median \emph{of the pulls
themselves}, which induces negative correlation among them, so the residual
correlation matrix is not the identity even when the underlying pulls are
independent. On independently generated series the estimate averages $8.57$ at
$m=9$ and $2.59$ at $m=3$, shortfalls of $4.8\%$ and $13.7\%$. The direction is
conservative --- it understates information and therefore errs toward rejecting a
batch --- but it means the effective strictness of
\texttt{min\_spectral\_effective\_pulls} exceeds its nominal value, most visibly
in small designs.

% =============================================================================
\section{Stage 6: temporal benchmarking}\label{sec:benchmark}

Long GT histories at daily frequency are not directly comparable to the weekly
or monthly series that GT returns for the same query over the full window. Let
$\bm{z}$ be the preliminary daily consensus, $A$ the aggregation matrix mapping
days to benchmark periods, and $\bm{b}$ the lower-frequency benchmark. The soft
benchmark solves
\begin{equation}
\min_{\bm{x}}\;
(\bm{x}-\bm{z})^{\top}Q(\bm{x}-\bm{z})
\;+\;
(A\bm{x}-\bm{b})^{\top}W(A\bm{x}-\bm{b}),
\label{eq:benchmark}
\end{equation}
where $Q = (\lambda_I+\lambda_R)I + \lambda_D D_1^{\top}D_1$ combines fidelity to
the daily consensus with a \emph{first}-difference smoothness penalty on the
correction path, and $W$ is the benchmark precision. The ridge $\lambda_R$ is a
separate configuration field whose only role is to keep $Q$ positive definite
when $\lambda_I=0$, so the two identity terms are one coefficient
mathematically. The default $\lambda_D/\lambda_I = 10$ is a strong smoothing
prior: it pushes corrections toward low-frequency adjustments and can
systematically damp daily shape near event days, which is where an attention
index is most informative. Calibrate and lock that ratio in a pilot, and report
sensitivity over $\lambda_D/\lambda_I \in \{1,10,100\}$. This is the
quadratic-programming form of the movement-preservation idea familiar from
\citet{Denton1971} and \citet{ChowLin1971}; treating $\bm{b}$ as noisy rather
than exact is appropriate here because the lower-frequency GT series is itself a
sampled quantity.

\begin{proposition}[Existence and uniqueness]\label{prop:bench}
If $Q \succ 0$ and $W \succeq 0$, the objective in \eqref{eq:benchmark} is
strictly convex and the solution is the unique
$\bm{x}^{*} = (Q + A^{\top}WA)^{-1}(Q\bm{z} + A^{\top}W\bm{b})$.
\end{proposition}

\begin{proof}
The Hessian is $2(Q + A^{\top}WA)$, which is positive definite because
$Q \succ 0$ and $A^{\top}WA \succeq 0$. Setting the gradient to zero gives the
stated normal equations, whose coefficient matrix is invertible.
\end{proof}

% =============================================================================
\section{Stage 7: reliability and the acceptance decision tree}\label{sec:decision}

Three reliability questions are kept separate because they can have different
answers for the same series: whether the zero/non-zero pattern reproduces
(detection, summarized by Fleiss' $\kappa$, \citealp{Fleiss1971}), whether
levels reproduce, and whether daily innovations reproduce.

The acceptance decision is a locked tree evaluated after estimation. Each rule
compares an observed diagnostic against a prespecified threshold and returns
pass, fail, or indeterminate. The batch status is
\begin{itemize}
\item \textbf{FAIL} if any mandatory rule fails;
\item \textbf{REVIEW} if no mandatory rule fails but at least one is
indeterminate;
\item \textbf{PASS} otherwise.
\end{itemize}

The rules cover protocol integrity, overlap-graph connectivity, the number of
numerically distinct pulls, the spectral effective pull count, the share of zero
observations, level and innovation generalizability, detection agreement,
concentration of the near-duplicate dependence graph, consensus convergence as
the last pull is added, and consistency with the lower-frequency benchmark.

\begin{remark}[Indeterminate is not a soft fail]
The detection rule is not identified when every observation falls in one
detection category, because $\kappa$ is undefined without variation. The
implementation reports the rule, marks it indeterminate, and---crucially---makes
it non-mandatory in exactly that case, so that a series with no zeros is not
penalized for a coefficient that cannot exist. Conflating ``cannot be computed''
with ``failed'' would systematically reject high-volume terms.
\end{remark}

\begin{remark}[What PASS does and does not mean]
PASS authorizes construction under the locked protocol. It does not establish
that the keyword has construct validity for the economic concept of interest, it
does not identify absolute search counts, and it licenses no causal
interpretation.
\end{remark}

% =============================================================================
\section{Stage 8: propagating measurement error downstream}\label{sec:eiv}

Let $W_t = X_t + u_t$ be the constructed index, used as a regressor for an
outcome $Y_t$ with controls. Under classical additive measurement error, the
moment-corrected estimator is
\begin{equation}
\hat{\beta} \;=\;
\frac{\bm{W}^{\top}M\bm{Y}}
{\bm{W}^{\top}M\bm{W} - \tr(M\Omega_u)},
\label{eq:rc-ols}
\end{equation}
where $M$ is the residual maker for the controls and $\Omega_u$ the measurement
error covariance supplied by the construction stage. Standard errors use a HAC
estimator \citep{NeweyWest1987}. The implementation raises an error when the
corrected denominator is non-positive, rather than returning an
explosive or sign-flipped coefficient; a non-positive denominator means the
estimated measurement error exceeds the observed regressor variation, which is a
statement about the data, not a numerical detail to be smoothed over.

As an alternative that does not require the classical structure to hold exactly,
SIMEX \citep{CookStefanski1994,Carroll2006} adds further known noise on a grid
of $\lambda$ values, records the resulting estimates, fits an extrapolation
function, and evaluates it at $\lambda = -1$.

% =============================================================================
\section{Controlled Monte Carlo evidence}\label{sec:monte-carlo}

The simulation generates a latent signal with realistic dynamics, applies
collection-day and stream effects, adds correlated retrieval error, splits the
history into independently normalized chunks, imposes rounding and zeros, and
optionally injects exact duplicates. Because the latent truth is known by
construction, estimators can be compared directly.

\subsection{Matching the information sets}

\texttt{PipelineResult.final\_series} returns the benchmarked series whenever a
benchmark is available. Comparing that series against cross-pull baselines that
never received the benchmark does not compare four estimators; it compares three
estimators against one estimator plus an extra block of low-frequency
information, and in simulation that block is generated from the latent truth
itself. Every estimator is therefore evaluated twice: on the calibrated pulls
alone, and after applying the \emph{same} weekly and monthly benchmark. The
single-pull row reports mean performance over all pulls rather than one fixed
pull, because a fixed column measures which pull happened to sort first, not
what a researcher who downloads once should expect.

\begin{table}[H]
\centering
\small
\caption{Twenty controlled replications (seeds 1200--1219). MCSE is the Monte
Carlo standard error of the mean RMSE. Lower RMSE is better.}
\label{tab:mc}
\begin{tabular}{llrrrrr}
\toprule
Estimator & Benchmark & RMSE & MCSE & Corr. & Innov.\ corr. & Peak recall\\
\midrule
Single pull        & no  & 7.3519 & 0.0573 & 0.9688 & 0.8878 & 0.8912\\
Cross-pull median  & no  & 4.1710 & 0.0819 & 0.9898 & 0.9556 & 0.9447\\
Simple mean        & no  & \textbf{3.7253} & 0.0877 & 0.9919 & \textbf{0.9747} & 0.9526\\
\RACER{}           & no  & 3.7421 & 0.0894 & 0.9918 & 0.9743 & 0.9526\\
\midrule
Single pull        & yes & 6.9881 & 0.0491 & 0.9717 & 0.8882 & 0.8947\\
Cross-pull median  & yes & 3.9584 & 0.0695 & 0.9908 & 0.9558 & 0.9474\\
Simple mean        & yes & \textbf{3.5058} & 0.0754 & \textbf{0.9928} & \textbf{0.9748} & \textbf{0.9553}\\
\RACER{}           & yes & 3.5226 & 0.0770 & 0.9927 & 0.9745 & 0.9526\\
\bottomrule
\end{tabular}
\end{table}

\begin{table}[H]
\centering
\small
\caption{Paired comparisons on RMSE. A positive difference favours estimator A;
$n=20$ paired replications.}
\label{tab:mcpaired}
\begin{tabular}{llrrrl}
\toprule
Estimator A & Estimator B & Difference & $t$ & $p$ & A wins\\
\midrule
\RACER{} & Single pull & $+3.6098$ & $48.31$ & $2.4\times10^{-21}$ & 20/20\\
\RACER{} & Cross-pull median & $+0.4289$ & $14.86$ & $6.5\times10^{-12}$ & 20/20\\
\RACER{} & Simple mean & $-0.0167$ & $-1.49$ & $0.152$ & 7/20\\
\RACER{} (bm.) & Simple mean (bm.) & $-0.0168$ & $-1.65$ & $0.116$ & 7/20\\
\RACER{} (bm.) & \RACER{} (no bm.) & $+0.2194$ & $15.49$ & $3.1\times10^{-12}$ & 20/20\\
Simple mean (bm.) & Simple mean (no bm.) & $+0.2195$ & $15.45$ & $3.3\times10^{-12}$ & 20/20\\
\bottomrule
\end{tabular}
\end{table}

Three conclusions follow, two of which support the framework and one of which
does not. First, cross-pull aggregation buys a great deal: at a matched
information set \RACER{} cuts RMSE by $1-3.7421/7.3519=49.1\%$ against a single
pull, improving in all twenty replications, and by a further $10.3\%$ against
the cross-pull median. Second, covariance-adjusted weighting shows no detectable
advantage over the simple mean, before benchmarking ($-0.0167$, $p=0.152$) or
after it ($-0.0168$, $p=0.116$). Third, temporal benchmarking contributes about
$0.22$ on its own, and contributes it almost identically to \RACER{} and to the
simple mean ($0.2194$ versus $0.2195$): the gain belongs to the benchmark, not
to the daily estimator it is paired with.

\begin{remark}[A testable prediction]\label{rem:prediction}
The second conclusion is a negative result, and its most natural explanation is
that dependence in this DGP is homogeneous: every pull shares its collection-day
noise with fixed weight $0.5$, so the residual covariance is close to
equicorrelated, the minimum-variance weights of an equicorrelated matrix are
equal weights, and GLS has nothing to exploit while still paying for estimating
$\Sig$.

If that explanation holds, GLS should win once dependence strength \emph{differs}
across pulls --- when some pulls are served from one cache and others are not.
That is a testable prediction, and \cref{sec:depexp} tests it directly.
\end{remark}

\subsection{A direct test of dependence heterogeneity}\label{sec:depexp}

\texttt{examples/run\_dependence\_experiment.py} places a subset of pulls behind
one shared cache disturbance. Membership cuts across collection days and streams,
so the extra dependence \emph{cannot} be absorbed by the design facets and is
visible only in the residual covariance --- which is what the GLS weights exist
for. At \texttt{cache\_cluster\_fraction} $=0$ the simulated data are byte-identical
to those behind \cref{tab:mc}, so the two sections are directly comparable.

\begin{table}[H]
\centering
\small
\caption{Four dependence scenarios, 20 replications each. Advantage is simple
mean minus GLS in RMSE, so positive favours GLS.}
\label{tab:depexp}
\begin{tabular}{lrrrrrrl}
\toprule
Scenario & GLS & Simple mean & Median & Advantage & MCSE & $p$ & GLS wins\\
\midrule
Homogeneous (baseline) & 3.7421 & 3.7253 & 4.1710 & $-0.0167$ & 0.0112 & $0.152$ & 7/20\\
1/3 of pulls, moderate & 3.7658 & 3.6801 & 4.1133 & $-0.0857$ & 0.0155 & $<0.0001$ & 2/20\\
1/3 of pulls, strong & 4.1684 & 4.1980 & 4.6486 & $+0.0296$ & 0.0439 & $0.508$ & 10/20\\
1/2 of pulls, strong & 5.1633 & 4.8802 & 5.5593 & $-0.2831$ & 0.0591 & $0.0001$ & 2/20\\
\bottomrule
\end{tabular}
\end{table}

\textbf{The prediction fails.} GLS wins no scenario significantly. The closest
case (one third of pulls, strong dependence) is $+0.0296$ with $p=0.51$ and
exactly 10 wins out of 20, which is a tie; in two scenarios GLS is
\emph{significantly worse}. Adding heterogeneity does not reverse the ranking.

Two diagnostics rule out the obvious misreadings of that result.

First, \textbf{the constraints are not the cause}. Relaxing them under the
strong one-third scenario (10 replications), the constrained default, the
uncapped variant, and the fully unrestricted solution all attain a mean RMSE of
$4.1752$ --- identical to the digit, because the unrestricted solution already
lies inside the feasible set. An empirical covariance gives $4.1785$ and a
diagonal one $4.2345$, against $4.2415$ for the simple mean over the same seeds.
The unrestricted estimator, which is BLUE in theory, does not win either.

Second, \textbf{the mechanism works}. The correlation between a pull's mean
residual correlation and the weight it receives is $-0.562$, negative in all ten
replications: GLS does find the pulls that share dependence and does move weight
away from them. Weight dispersion in \cref{tab:depexp} rises monotonically with
heterogeneity ($0.0220\to0.0409$) and the spectral effective pull count falls
monotonically ($4.93\to2.95$); both correctly register the added dependence.

What remains is \textbf{estimation error}. $\Sig$ must be estimated from 9 pulls
over 366 historical dates, and the perturbation that its error induces in the
weights exceeds the efficiency that exploiting the true dependence structure can
buy. This matches a long-standing finding in the portfolio literature, where
sample-covariance minimum-variance portfolios fail to beat equal weighting in
small samples; Ledoit--Wolf shrinkage \citep{LedoitWolf2004} was proposed to
mitigate exactly this, and is already the default here.

\begin{remark}[What the negative result changes, and what it does not]
It does not change the theory. With $\Sig$ \emph{known}, the GLS weights remain
minimum-variance among linear unbiased estimators (\cref{prop:blue}); that
derivation depends on no simulation. It changes the practical recommendation: at
the design sizes this paper recommends (21--42 pulls, a single year), no
detectable gain in point accuracy should be expected from the covariance
adjustment, and a simple mean is an equally defensible primary choice.

It also does not change the \emph{diagnostic} value of the adjustment. The
spectral effective pull count fell from 4.93 to 2.95, correctly telling the
researcher that nine pulls are now worth about three. A simple mean cannot supply
that judgment, nor the near-duplicate dependence structure, the per-day
conditional standard errors, or measurement uncertainty in a form the downstream
EIV correction can consume. RMSE cannot answer whether a batch is usable at all,
and that is what this framework is mainly for.
\end{remark}

\begin{remark}[Reading the bias column]
Mean bias for the single pull and for the simple mean sits at the $10^{-15}$
level. That is an identity, not an estimator property: every complete pull is
normalized to baseline mean 100 and the truth is defined on the same baseline,
so their average cannot be biased. \RACER{} now lands at machine precision as
well ($+3.2\times10^{-17}$). The cross-pull median's $-0.1464$ is real: a daily
median sits systematically below the mean under the right-skewed event spikes in
this DGP.
\end{remark}

% =============================================================================
\section{Scope of the statistical guarantees}\label{sec:scope}

We state precisely what is and is not established.

\begin{enumerate}
\item Under \cref{as:protocol,as:inclusion}, the stratified design-cell average
is design-unbiased for the prespecified retrieval-mixture expectation
(\cref{prop:design}). This holds \emph{without} any independence assumption
across pulls.
\item Under \cref{as:overlap,as:connected,as:edge-mean}, the relative log-scale
parameters are identified and the graph WLS estimator is conditionally unbiased
and efficient among linear estimators built on the edge estimates
(\cref{prop:ident,prop:gwls}).
\item Under \cref{as:common} with $\Sig$ known or consistently estimated, the
unconstrained GLS weights are minimum-variance among linear unbiased estimators
(\cref{prop:blue}). The constrained default is a stabilized convex combination
and does not inherit this optimality.
\item Given a correct lower-frequency benchmark and an identified constraint
set, the benchmark estimator is the unique minimizer of an explicit quadratic
objective (\cref{prop:bench}).
\item Under classical additive measurement error with a consistent estimate of
$\Omega_u$, \eqref{eq:rc-ols} corrects attenuation bias.
\end{enumerate}

\noindent None of these results establishes unbiasedness with respect to
Google's absolute search counts, and none is claimed to hold unconditionally.

\subsection{Known limitations}

\begin{itemize}
\item The variance components use balanced method-of-moments estimation. REML
for unbalanced designs is a planned extension.
\item Consensus confidence intervals are conditional approximations. A
full-pipeline moving-block bootstrap---resampling historical blocks and
re-estimating edges, scales, duplicates, covariance, weights, and the benchmark
before re-estimating the financial model---is the appropriate procedure for
confirmatory inference.
\item Benchmark noise, scale-calibration uncertainty, and tuning-parameter
uncertainty are not yet fully propagated into the reported intervals.
\item Zero values may mix genuine zero search, privacy thresholds, sampling, and
rounding. Version \racerversion{} retains zeros and isolates detection
reliability, but does not implement an interval-censored latent-count
likelihood. For high zero-share terms, moving to weekly frequency or to a
broader Topic is usually more defensible than manufacturing daily values.
\item Anchor-bank calibration across different queries is documented as an
extension but is not implemented in version \racerversion{}.
\item \textbf{Real-data validation covers the calibration stage only.}
\Cref{sec:realdata} uses a single pull, so every mechanism from Stage 2 onward
--- consensus weighting, duplicate diagnosis, variance components, reliability
coefficients, benchmarking --- is \emph{entirely untested} on real Google Trends
data and rests on simulation alone. \Cref{tab:chunknoise} has already shown the
simulator overstating chunk noise by about a factor of 4.5 in the one stage where
a comparison was possible, which makes it an open and consequential question
whether its other parameters are distorted the same way: if cross-pull retrieval
noise, day effects, or stream effects are also overstated, the relative
performance figures in \cref{sec:monte-carlo} need re-examination too. We
\textbf{do not claim} that simulation evidence extrapolates to cross-pull
behaviour on real data. Closing the gap requires repeated collection across
collection days and stream environments; under the recommended protocol, three
streams by seven collection days is 21 pulls and 168 downloads, which is a
month of calendar time.
\item The package is deliberately ingestion-first and embeds no unofficial
scraper, because endpoints, rate limits, authentication, and terms of service
change. Collectors need only emit the documented long-format schema.
\item \textbf{Retrieval rate limiting is a binding constraint, and it undermines
the design rather than merely slowing it.} The confirmatory protocol recommended
above --- from 2010, 180-day windows, 15--30 day steps, 21--42 pulls --- implies a
request budget that was never stated. At a 15-day step a single pull spans 393
chunks, and each chunk costs two requests (explore for a token, then
widgetdata), so one pull is 786 requests, 21 pulls is $16{,}506$, and 42 pulls is
$33{,}012$. At a conservative 20-second spacing those are roughly 92 and 183
hours of continuous retrieval; even one pull takes about 4.4 hours.

Google Trends rate limits this aggressively. In one test the explore endpoint
returned 200 while the \emph{second} widgetdata request returned HTTP 429, and a
retry 30 seconds later returned 429 again. A single observation cannot quantify
Google's policy, which may vary with address, time of day, account state, and
request pattern, but it is enough to establish that acquisition feasibility is
not a detail that can be left to the collector.

The limit is not confined to programmatic access. Loading the Trends web page in
an ordinary signed-in browser during the same period, by hand, produced a 429 on
the front end's \emph{own} \texttt{widgetdata} request while the chart rendered
normally from cache. That has an immediate operational consequence:
\textbf{the Trends front end masks failed requests with cached content, so a user
cannot tell from the screen whether the series shown was actually retrieved in
this session}. Any protocol that relies on manual export must record the response
status and \texttt{raw\_response\_hash} independently rather than treating a
rendered chart as evidence of successful retrieval. This sharpens the reason the
adapter contract insists on preserving the raw response.

The damage is not the waiting. \Cref{sec:design} defines \texttt{collection\_day}
as the ordinal level of one complete historical reconstruction, and both the $D$
facet of the G-study and $\sigma^2_{TD}$ rest on that level. If one complete
reconstruction takes 4.4 hours and is repeatedly interrupted by throttling, which
collection day a pull belongs to stops being well defined; a pull that crosses
midnight additionally straddles the cache reset described by
\citet{Djorno2026} (\cref{rem:cache}), so its chunks come from two different
sampling batches. That is not measurement error, it is the failure of a facet
definition.

The remedy is to size the window against the limit rather than the other way
round. \texttt{examples/collect\_google\_trends.py --preview} prints the request
budget before any data are collected, and a single year with a 180-day window and
a 30-day step needs 8 chunks and 16 requests per pull, completing in minutes ---
short enough to keep the collection day well defined. A 16.6-year window at a
15-day step is not currently feasible, and earlier versions of this paper did not
say so.
\end{itemize}

% =============================================================================
\section{Conclusion}

Google Trends is a measurement instrument with a documented tendency to return
different numbers for the same question. Treating its output as an
error-free observation is not a conservative simplification; it silently imports
scale error, dependence, and attenuation into whatever model consumes it.

\RACER{} takes the opposite stance. It fixes the retrieval design before any
data are seen, calibrates all chunks simultaneously on the overlap graph,
measures rather than assumes the independence of repeated retrievals,
decomposes the remaining variation across the facets of the design, and carries
the resulting uncertainty forward into downstream estimation. What it delivers
is not certainty about Google's internal counts, which is unavailable, but a
reproducible, auditable, and explicitly conditional estimate of a well-defined
target, together with a locked rule for deciding whether a given batch of data
is fit to use at all.

The controlled evidence is deliberately mixed. Under this data-generating
process, cross-pull aggregation and temporal benchmarking both deliver large and
robust reductions in reconstruction error, while covariance-adjusted weighting
does not beat a simple mean. We report the comparison that does not favour us
alongside the ones that do, because that is the same discipline the framework
asks of its users: a measurement procedure that can only produce flattering
numbers cannot tell a researcher when to stop trusting it.

% =============================================================================
\begin{thebibliography}{99}
\small

\bibitem[Brennan(2001)]{Brennan2001}
Brennan, R.~L. (2001). \textit{Generalizability Theory}. Springer, New York.

\bibitem[Carroll et al.(2006)]{Carroll2006}
Carroll, R.~J., Ruppert, D., Stefanski, L.~A., and Crainiceanu, C.~M. (2006).
\textit{Measurement Error in Nonlinear Models: A Modern Perspective},
2nd edition. Chapman and Hall/CRC, Boca Raton.

\bibitem[Choi and Varian(2012)]{ChoiVarian2012}
Choi, H. and Varian, H. (2012). Predicting the present with Google Trends.
\textit{Economic Record}, 88(s1), 2--9.

\bibitem[Chow and Lin(1971)]{ChowLin1971}
Chow, G.~C. and Lin, A. (1971). Best linear unbiased interpolation,
distribution, and extrapolation of time series by related series.
\textit{Review of Economics and Statistics}, 53(4), 372--375.

\bibitem[Cook and Stefanski(1994)]{CookStefanski1994}
Cook, J.~R. and Stefanski, L.~A. (1994). Simulation-extrapolation estimation in
parametric measurement error models. \textit{Journal of the American
Statistical Association}, 89(428), 1314--1328.

\bibitem[Cronbach et al.(1972)]{Cronbach1972}
Cronbach, L.~J., Gleser, G.~C., Nanda, H., and Rajaratnam, N. (1972).
\textit{The Dependability of Behavioral Measurements: Theory of
Generalizability for Scores and Profiles}. Wiley, New York.

\bibitem[Da et al.(2011)]{DaEG2011}
Da, Z., Engelberg, J., and Gao, P. (2011). In search of attention.
\textit{Journal of Finance}, 66(5), 1461--1499.

\bibitem[Denton(1971)]{Denton1971}
Denton, F.~T. (1971). Adjustment of monthly or quarterly series to annual
totals: An approach based on quadratic minimization. \textit{Journal of the
American Statistical Association}, 66(333), 99--102.

\bibitem[Djorno et al.(2026)]{Djorno2026}
Djorno, C., Santillana, M., and Yang, S. (2026). Restoring the forecasting power
of Google Trends with statistical preprocessing. \textit{International Journal
of Forecasting}, 42(3), 1104--1122. Preprint: arXiv:2504.07032.

\bibitem[Eichenauer et al.(2022)]{EichenauerIMS2022}
Eichenauer, V.~Z., Indergand, R., Mart\'{i}nez, I.~Z., and Sax, C. (2022).
Obtaining consistent time series from Google Trends.
\textit{Economic Inquiry}, 60(2), 694--705.

\bibitem[Bleher and Dimpfl(2022)]{BleherDimpfl2022}
Bleher, J. and Dimpfl, T. (2022). Knitting multi-annual high-frequency Google
Trends to predict inflation and consumption. \textit{Econometrics and
Statistics}, 24, 1--26.

\bibitem[Cebri\'{a}n and Domenech(2024)]{CebrianDomenech2024}
Cebri\'{a}n, E. and Domenech, J. (2024). Addressing Google Trends
inconsistencies. \textit{Technological Forecasting and Social Change}, 202,
123318.

\bibitem[Da et al.(2015)]{DaEG2015}
Da, Z., Engelberg, J., and Gao, P. (2015). The sum of all FEARS: Investor
sentiment and asset prices. \textit{The Review of Financial Studies}, 28(1),
1--32.

\bibitem[Fleiss(1971)]{Fleiss1971}
Fleiss, J.~L. (1971). Measuring nominal scale agreement among many raters.
\textit{Psychological Bulletin}, 76(5), 378--382.

\bibitem[H\"{o}lzl et al.(2025)]{Holzl2025}
H\"{o}lzl, J., Keusch, F., and Sajons, C. (2025). The (mis)use of Google Trends
data in the social sciences. \textit{Social Science Research}, 126, 103099.

\bibitem[Huber(1964)]{Huber1964}
Huber, P.~J. (1964). Robust estimation of a location parameter.
\textit{Annals of Mathematical Statistics}, 35(1), 73--101.

\bibitem[Kish(1965)]{Kish1965}
Kish, L. (1965). \textit{Survey Sampling}. Wiley, New York.

\bibitem[Ledoit and Wolf(2004)]{LedoitWolf2004}
Ledoit, O. and Wolf, M. (2004). A well-conditioned estimator for
large-dimensional covariance matrices. \textit{Journal of Multivariate
Analysis}, 88(2), 365--411.

\bibitem[Medeiros and Pires(2021)]{MedeirosPires2021}
Medeiros, M.~C. and Pires, H.~F. (2021). The proper use of Google Trends in
forecasting models. arXiv:2104.03065.

\bibitem[Newey and West(1987)]{NeweyWest1987}
Newey, W.~K. and West, K.~D. (1987). A simple, positive semi-definite,
heteroskedasticity and autocorrelation consistent covariance matrix.
\textit{Econometrica}, 55(3), 703--708.

\bibitem[Neumann et al.(2023)]{Neumann2023}
Neumann, K., Mason, S.~M., Farkas, K., Santaularia, N.~J., Ahern, J., and
Riddell, C.~A. (2023). Harnessing Google Health Trends data for epidemiologic
research. \textit{American Journal of Epidemiology}, 192(3), 430--437.

\bibitem[Pagan(1984)]{Pagan1984}
Pagan, A.~R. (1984). Econometric issues in the analysis of regressions with
generated regressors. \textit{International Economic Review}, 25(1), 221--247.

\bibitem[Google Search Central(2025)]{GoogleAPI2025}
Google Search Central (2025). Introducing the Google Trends API (alpha): A new
way to access Trends data. Published 24 July 2025.

\bibitem[Google Trends Help(n.d.)]{GoogleFAQ}
Google Trends Help (n.d.). FAQ about Google Trends data.

\bibitem[Preis et al.(2013)]{PreisMS2013}
Preis, T., Moat, H.~S., and Stanley, H.~E. (2013). Quantifying trading behavior
in financial markets using Google Trends. \textit{Scientific Reports}, 3, 1684.

\bibitem[West(2020)]{West2020}
West, R. (2020). Calibration of Google Trends time series. In
\textit{Proceedings of the 29th ACM International Conference on Information and
Knowledge Management}, 2257--2260.

\end{thebibliography}

\end{document}
